\scanpage{21}
sont isomorphes respectivement à
\[
\operatorname{Coker}\big(H^0(A'_0/A_0,G_m)\to H^1(IK(A'/A))\big)
\]
et à
\[
\operatorname{Coker}\big(H^0(A'_0/A_0,G_a)\to H^1(IK(A'/A))\big)
\]

Or on a un diagramme commutatif d'applications
\diagram{\begin{tikzcd}[ampersand replacement=\&]
B^* \arrow[r] \arrow[d,"i"] \& C_0^* \arrow[r] \arrow[d,"j"] \& H^1(IK(A'/A)) \arrow[d,"\vert\,\mathrm{id}"] \\
B \arrow[r] \& C_0 \arrow[r] \& H^1(IK(A'/A))
\end{tikzcd}}
où $i,j$ sont les \unsure{immersions|inclusions} canoniques (elles
sont pas multiplicatives, mais cela est sans importance).
Le résultat résulte de cela, et de l'hypothèse
\[
C_0=C_0^*+\operatorname{Im}(B\to C_0).
\]

\uline{Corollaire} 1 Supposons que
\[
H^1(A'_0/A_0,G_a)=0\qquad H^1(A'_0/A_0,G_m)=0
\]
\deleted{\unsure{Pour qu'on ait|Pour qu'il existe}} \inserted{Alors on a} un \deleted{bijection} isomorphisme can.
\[
H^1(A'/A,G_m)\simeq H^1(A'/A,G_a)
\]
donc les relations $H^1(A'/A,G_m)=0$ et
$H^1(A'/A\unsure{,}G_a)=0$ sont équivalentes.
% archive sheet number: 20
